%% kapa: exam=1 round=26 year=2015 period=1 form=A subject=civil_law subject_ko=민법 seq=31 number=31
%% source: Q-Net 감정평가사 기출문제 (cst00309) artlSeq=5138001 file=extracted/1008/5138001/제26회_감정평가사 1차_기출/제26회_감정평가사 1차_기출/1교시_감평_A형.pdf
%% conversion: kordoc 4.24.51 -> tools/gen1.py (auto); review=proofread-sonnet; flags=none
\begin{question}{31}
주위토지통행권에 관한 설명으로 옳은 것은? (다툼이 있으면 판례에 따름)
\begin{choices}
  \item 통행권자가 통행지 소유자에게 손해보상의 지급을 게을리 하면 통행권이 소멸한다.
  \item 주위토지통행권의 범위는 현재의 토지의 용법은 물론 장래의 이용 상황도 미리 대비하여 정해야 한다.
  \item 통행권자가 통행지를 배타적으로 점유하는 경우 통행지 소유자는 통행지의 인도를 청구할 수 있다.
  \item 주위토지통행권이 인정되는 경우 통행지 소유자는 원칙적으로 통로개설 등 적극적인 작위의무를 부담한다.
  \item 동일인 소유의 토지의 일부가 양도되어 공로에 통하지 못하는 토지가 생긴 경우, 포위된 토지를 위한 통행권은 일부 양도 전의 양도인 소유의 종전 토지뿐만 아니라 다른 사람 소유의 토지에 대하여도 인정된다.
\end{choices}
\end{question}
\answer{3}
